\documentclass[11pt]{article}
\usepackage[utf8]{inputenc}	% Para caracteres en español
\usepackage{amsmath,amsthm,amsfonts,amssymb,amscd}
\usepackage{multirow,booktabs}
\usepackage[table]{xcolor}
\usepackage{fullpage}
\usepackage{lastpage}
\usepackage{enumitem}
\usepackage{fancyhdr}
\usepackage{mathrsfs}
\usepackage{algpseudocode}
\usepackage{wrapfig}
\usepackage{xcolor}
\usepackage{setspace}
\usepackage{calc}
\usepackage{multicol}
\usepackage{cancel}
\usepackage{tikz}
\usepackage[retainorgcmds]{IEEEtrantools}
\usepackage[margin=3cm]{geometry}
\usepackage{amsmath}
\newlength{\tabcont}
\setlength{\parindent}{0.0in}
\setlength{\parskip}{0.05in}
\usepackage{empheq}
\usepackage{framed}
\usepackage{multicol}
\usepackage{amsmath}
\usepackage{setspace}
\usepackage{graphicx}
\usepackage{caption}
\usepackage{subcaption}
\usepackage{graphbox}
\graphicspath{ {figures/} }
\usepackage[most]{tcolorbox}
\usepackage{xcolor}
\colorlet{shadecolor}{orange!15}
\parindent 0in
\parskip 12pt
\geometry{margin=1in, headsep=0.25in}
\theoremstyle{definition}
\newtheorem{defn}{Definition}
\newtheorem{reg}{Rule}
\newtheorem{exer}{Exercise}
\newtheorem{note}{Note}
\newcommand{\indep}{\perp\!\!\!\!\perp} 
\begin{document}
\title{Report}

\thispagestyle{empty}

\begin{center}
{\LARGE \bf Large Sample Theory Exercises}\\
\end{center}
\tableofcontents
\newpage
\section{Convergence of Random Variables}
1. (2022 Final) Prove or disprove (with a counter example): $X_n \rightarrow X$, completely implies $X_n^2 \rightarrow X^2$, almost surely.

2. Suppose we have the following representation
$$
T_n=Z_n+R_n Q_n
$$
where $E\left(\left|R_n\right|\right) \rightarrow 0, \quad Q_n \xrightarrow{d} Q$, for some random variable $Q$, and $Z_n \xrightarrow{d} N(0,1)$. What will be the limiting distribution of $T_n^2$ ?

3. Suppose $X_1, \cdots, X_n, \cdots$ be i.i.d. Uniform $(-1,1)$. Show from definition that $\bar{X}_n \rightarrow 0$ in $L_2$, where $\bar{X}_n=\frac{1}{n} \sum_{i=1}^n X_i$.

4. Suppose $X_n=n^2 Y_n$, where $Y_n$ are defined on the same probability space such that $Y_n$ has a Bernoulli distribution with probability of success $1 / n^2$.

(i) Show that $X_n$ converges to the 0 random variable, almost surely, as $n \rightarrow \infty$.

(ii) Does $X_n$ converge to the 0 random variable in $L_1$ ?

5. Suppose $X_n$ are defined on the same probability space and that

$$
\sum_{n=1}^{\infty} E\left(\left|X_n\right|^p\right)<\infty, \text { for some } p>0
$$
Show that

(a) $X_n \rightarrow 0$, completely.

(b) Is it true that $X_n \rightarrow 0$, in probability?

(c) Is it true that $X_n \rightarrow 0$, in $L_1$, if $p \geq 1$ ?


6. Prove or give counter examples.

(1). If $X_n \rightarrow c$, in probability, then $X_n \rightarrow c$ in $L_1$.

(2). If $X_n \rightarrow X$ in $L_1, Y_n \rightarrow c$ in probability then $X_n+Y_n \rightarrow X+c$ in $L_1$.

(3). If $X_n \rightarrow X$ in $L_1, Y_n \rightarrow c$ in probability and $\left|Y_n\right| \leq 1$, for all $n$ with probability one, then $X_n+Y_n \rightarrow X+c$ in $L_1$.

(4). If $E\left(X_n-Y_n\right)^2=O\left(n^{-2}\right)$ and $\sqrt{n}\left(X_n-c\right) \xrightarrow{d} N(0,5)$, then $Y_n \sim A N(c, 5 / n)$.

7. Prove or disprove (with counter example):

(a) $X_n \xrightarrow{P} X$ implies $X_n \rightarrow X$ completely.

(b) $X_n \sim A N\left(\mu_n, 1 / n\right)$ implies $X_n \sim A N(0,1 / n)$, provided $\mu_n \rightarrow 0$.

8. Suppose $X_i$ are iid with $E\left(X_i\right)=\mu, \operatorname{Var}\left(X_i\right)=1$.

(i) What is the MSE of $\widehat{\mu}=\bar{X}$ ?

(ii) Does $\widehat{\mu}$ converge to $\mu$ in $L_1$ ? Give arguments for your answer.

(iii) What is $\operatorname{Var}_{J K}(\widehat{\mu})$ in this case?

9. Ture or false? If you claim it to be true, give a short proof; otherwise, give a counter example.

(a) If $\mathrm{E}\left(X_n\right)<\pi, \forall n$, then $X_n$ are uniformly integrable.

(b) If $E\left(|X|+\left|X_n\right|\right)<\infty, \forall n$, and $E\left(\left|X-X_n\right|\right) \rightarrow 0$, as $n \rightarrow \infty$, then $X_n$ are uniformly integrable.

10. TRUE or FALSE? Give short proofs or counter examples.

(a) $\sum_{n=1}^{\infty} P\left(\left|X_n\right|>n^{-1}\right)<\infty$ and $\left|X_n\right| \leq 1$ with probability 1 together imply $E\left|X_n\right| \rightarrow 0$.

(b) $\sum_{n=1}^{\infty} P\left(\left|X_n\right|>n^{-1}\right)<\infty$ implies $E\left|X_n\right| \rightarrow 0$.

(c) $\hat{\theta}_n \sim A N\left(\theta, \sigma^2 / n\right)$ implies $E\left(\hat{\theta}_n\right)=\theta$ and $\operatorname{Var}\left(\hat{\theta}_n\right)=\sigma^2 / n$.

11. Suppose $X_i \sim \operatorname{Bin}\left(1, p_i\right), i=1, \ldots, n$ are independent. Let $\widehat{p}=\frac{1}{n} \sum_{i=1}^n X_i$ and $\bar{p}=\frac{1}{n} \sum_{i=1}^n p_i$.

(a) Show that $\widehat{p}-\bar{p} \xrightarrow{P} 0$ as $n \rightarrow \infty$. 

(b) Show that if $\lim _{i \rightarrow \infty} p_i=p$, then $\widehat{p} \xrightarrow{P} p$, as $n \rightarrow \infty$.

(c) Under the condition in (b), show that $\widehat{p} \xrightarrow{L_2} p$ also.

12. For each of the following, state if the statement is true or false. If true, give a short proof, and if false, provide a counter example. You can use any result covered in class - just state the result you are using.

(a) If $\sup _n \mathbb{E}\left(X_n^2\right)<\infty$ and $X_n \xrightarrow{P} 0$, then $\mathbb{E}\left(X_n\right) \rightarrow 0$.

(b) A jackknife variance estimator is always consistent, i.e., the ratio of the jackknife variance estimator and the true variance converges almost surely to 1.

(c) If $X_n^2 \xrightarrow{P} 0$ and $-1 \leq X_n<2$ for all $n$, then $X_n \rightarrow 0$, almost surely.

13. True or False? If true give a short proof; if false, give a counter example.

(a) $X_n \xrightarrow{L_1} X$ and $Y_n \xrightarrow{L_1} Y$ together implies $X_n Y_n \xrightarrow{L_1} X Y$.

(b) $X_n \xrightarrow{L_2} X$ and $Y_n \xrightarrow{L_2} Y$ together implies $X_n Y_n \xrightarrow{L_1} X Y$.

(c) $U_n$ and $V_n$ are the $U$ and $V$ statistics, respectively, based on the same kernel $h$. Then
$$
n U_n V_n \xrightarrow{d} c \chi_1^2
$$
for some constant $c$.

14. Suppose $X_1, \cdots, X_n$ are i.i.d. with $0<\sigma^2=\operatorname{Var}\left(X_1\right)<\infty$. Suppose $\pi_i$ 's are i.i.d. taking values -1 and 1 with equal probabilities. Assume furthermore that the $X$ and the $\pi$ are independent. Show that $n^{-1} \sum_{i=1}^n \pi_i X_i \rightarrow 0$, in probability.

15. Suppose $X_n \xrightarrow{d} X, Y_n \xrightarrow{d} Y$. Assume furthermore, $X_n$ and $Y_n$ are independent for each $n$. Prove that $X_n+Y_n \xrightarrow{d} W$, for some $W$. Identify/characterize $W$.

16. Suppose $X_n=X+\frac{R_n}{n}$, for $n \geq 1$, where $R_n$ has a distribution function

$$
F_{R_n}(x)= \begin{cases}0, & \text { for } x<0 \\ (1-c) x, & \text { for } 0 \leq x<1 \\ 1-\frac{c}{x^2}, & \text { for } x \geq 1\end{cases}
$$

for some $0<c<1$.
Show that $X_n \rightarrow X$, almost surely.

17. True or False? Provide arguments supporting your answer:
$$
\text { If } \sum_{n=1}^{\infty} \mathbb{E}\left(\hat{\theta}_n-\theta\right)^2<\infty, \quad \text { then } \quad \operatorname{Bias}\left(\hat{\theta}_n\right) \rightarrow 0 .
$$

18. Prove or disprove (with a counterexample):

(a) If $\operatorname{Bias}(\hat{\theta}) \rightarrow 0$ and $\operatorname{Var}(\hat{\theta}) \rightarrow 0$, then $\hat{\theta}$ is weakly consistent; i.e. $\hat{\theta} \xrightarrow{P} \theta$.

(b) If $\operatorname{Bias}(\hat{\theta}) \rightarrow 0$ and
$$
\sum_{n=1}^{\infty}\{\operatorname{Bias}(\hat{\theta})+\operatorname{Var}(\hat{\theta})\}<\infty,
$$
where $n$ is the sample size, then $\hat{\theta}$ is strongly consistent; i.e. $\hat{\theta} \xrightarrow{\text { a.s. }} \theta$.

19. (2026 Exam 1) Consider i.i.d. data $X_1, \cdots, X_n$ from a parametric model $P_\theta: \theta \in \Re$. Suppose we have a regular (don't worry about the exact meaning of regular) estimator of $\theta$ with influence function $I F_{\widehat{\theta}}\left(X_i\right)=g\left(X_i, \theta\right)$, say.

(a) What does it mean (as stated in class)?

(b) Suppose $E\left(\widehat{\theta}^2\right)<\infty$, for all $n$, and $\sup _n E\left(r_n^{2+\delta}\right)<\infty$, for some $\delta>0$, where
$$
r_n=\sqrt{n}(\hat{\theta}-\theta)-\frac{1}{\sqrt{n}} \sum_{i=1}^n g^c\left(X_i, \theta\right) .
$$

(i) Show that $\operatorname{Var}\left(r_n\right) \rightarrow 0$, and

(ii) $\frac{\operatorname{Var}(\widehat{\theta})}{\frac{1}{n^2} \sum_{i=1}^n\left\{g^c\left(X_i, \theta\right)\right\}^2} \xrightarrow{p} 1$, and hence $\operatorname{Var}(\widehat{\theta})$ can be estimated by $\frac{1}{n^2} \sum_{i=1}^n\left\{g^c\left(X_i, \widehat{\theta}\right)\right\}^2$.

20. (2026 HW2) Ture or False? Give arguments or a counter example for your answer.

Suppose $X_1, \cdots, X_n$ are i.i.d. such that $E\left(X_1\right)=1, \mathrm{E}\left(X_1^2\right)=3$ and $E(X_1^4)<\infty$

Then $\frac{X_1+X_2+\cdots+X_{\mathrm{n}}}{n} \xrightarrow{c} 1$.

21. (2026 HW2) Ture or False? Give arguments or a counter example for your answer.

Suppose $X_n \geq 0$ with probability one, $\forall n, X_n \xrightarrow{P} 0$, and $\limsup _{n \rightarrow \infty} E\left(X_n^2\right)<\infty$. Then $E\left(X_n\right) \rightarrow 0$, as $n \rightarrow \infty$.

22. (2026 HW2) Suppose $X_i$ are i.i.d. with $\mu=E\left(X_i\right), \sigma^2:=\operatorname{var}\left(X_i\right) \in(0, \infty)$.

Let $\bar{X}_n=\left(X_1+X_2+\cdots+X_n\right) / n$,
$$
\sigma_n^2=\frac{1}{n} \sum_{i=1}^n\left(X_i-\bar{X}_n\right)^2 \text { and } s_n^2=\frac{1}{n-1} \sum_{i=1}^n\left(X_i-\bar{X}_n\right)^2 .
$$
Show that
$$
\sqrt{n}\left(\bar{X}_n-\mu\right) / s_n=\sqrt{n}\left(\bar{X}_n-\mu\right) / \sigma_n+\mathrm{o}_p(1) .
$$

23. (2026 Final) Suppose $X_n \xrightarrow{d} X, Y_n \xrightarrow{d} Y$. Furthermore, suppose, $X_n$ and $Y_n$ are independent for each $n$.

(a) Show that $X_n+Y_n \xrightarrow{d} X^*+Y^*$, where $X^* \stackrel{d}{=} X, Y^* \stackrel{d}{=} Y$, and $X^*$ and $Y^*$ are independent. (Hint: You may use a characteristic function argument)

(b) In particular if $X^* \sim N\left(a, \sigma^2\right), Y^* \sim N\left(b, \tau^2\right)$ then $X_n+Y_n \xrightarrow{d} N\left(a+b, \sigma^2+\tau^2\right)$.

24. (2025 Final) Suppose $X_i$ are iid with a symmetric and continuous distribution function $F$. Let, for any $x \in \Re$,
$$
F_n(x)=\frac{1}{n} \sum_{i=1}^n I\left(X_i \leq x\right)
$$
be the proportion of data values that are less than or equal to $x$. Note that $F_n$ is also called the empirical cdf.

(i) Ture or False? $F_n(0) \rightarrow 0.5$, almost surely. Give arguments in support of your answer.

(ii) Ture or False? $F_n(0) \rightarrow 0.5$, in $L_1$. Give arguments in support of your answer.

(iii) Establish the joint asymptotic normality of $\left(F_n(0), \widehat{\xi}_{1 / 2}\right)$, where $\widehat{\xi}_{1 / 2}$ is the sample median.

State your assumptions on $F$ for the validity of your answer.

\newpage

\section{Asymptotic Normality and MLE}
1. Consider $X_1, \cdots, X_n$ i.i.d. from a distribution $F$ and $Y_1, \cdots, Y_n$ i.i.d. from a distribution $G$. Furthermore, the $X$ s are independent of the $Y$ s. Suppose
$$
s_1^2=\frac{1}{n-1} \sum_{i=1}^n\left(X_i-\bar{X}\right)^2 ; s_2^2=\frac{1}{n-1} \sum_{i=1}^n\left(Y_i-\bar{Y}\right)^2 .
$$
Establish the asymptotic normality of $s_1^2 / s_2^2$.

2. Consider a location family $f_\theta(x)=f(x-\theta), \theta \in \Re$, where $f$ is a given density which is symmetric around 0 . Consider estimating $\theta$ by solving an estimating equation of the form

$$
\sum_{i=1}^n \mathrm{w}_i \psi\left(X_i-\theta\right)=0
$$

where $\mathrm{w}_i$ are known weights.

(a) What should we assume about $\psi$ so that the above estimating equation will lead to an unbiased estimator? 

(b) Find an expression for an asymptotic confidence interval for $\theta$ based on such an estimator $\widehat{\theta}$.

3. Consider i.i.d. data $X_1, \cdots, X_n$ and the following robust measure of the coefficient of variation
$$
\widehat{\theta}=\frac{\widehat{\tau}}{M},
$$
where
$$
\widehat{\tau}=\frac{1}{n^2} \sum_{i=1}^n \sum_{j=1}^n\left|X_i-X_j\right|
$$
and $M$ is the sample median.
Construct a variance estimator for $\widehat{\theta}$ using each of the following methods

(a) the large sample theory techniques (influence function delta method etc)

(b) the jackknife

(c) the bootstrap

Which one is not likely to work?

4. If $\widehat{\theta}$ is an asymptotically normal estimator of $\theta$ with influence function $A_i= A\left(X_i, \theta\right)$, and $g$ is a smooth function, what will be the influence functions of 

(i) $g(\widehat{\theta})$ 

(ii) $\widehat{\theta} \cdot g(\widehat{\theta})$

5. Let $X_1, \cdots, X_n$ be i.i.d., with density $f_\theta(x)=\theta^{-1} \exp \{-x / \theta\}, x>0$. Under what condition on the function $\gamma$, the following estimating function $\psi(X, \theta)=\gamma(X / \theta)-1$ be unbiased for zero? Give an example of such a $\gamma$. Write down the asymptotic distribution of $\widehat{\theta}$ obtained from the estimating equation
$$
\sum_{i=1}^n\{\gamma(X / \theta)-1\}=0
$$

for a general $\gamma$ satisfying the condition.

6. Suppose $X_1, \cdots, X_n$ are i.i.d. $N\left(\mu, \sigma^2\right)$. Find the asymptotic distribution of
$$
\widehat{\theta}=\frac{\bar{X}+\widehat{\xi}_{1 / 2}}{2},
$$
where $\bar{X}$ is the sample mean and $\widehat{\xi}_{1 / 2}$ is the sample median.

7. Consider a linear regression model $Y_i=\beta^T X_i+\epsilon_i, 1 \leq i \leq n, \beta \in \Re^p$, where $\epsilon_i \sim i . i . d . N(0,1)$, and are independent of the $X_i$ (in case the covariates are random). Consider estimating the regression parameters $\beta$ by estimating equations of the form
$$
\sum_{i=1}^n \psi\left(Y_i-\beta^T X_i\right) X_i=0
$$
Under what condition on the $\psi$, will you get a set of unbiased estimating equations? For such a $\psi$ find the asymptotic distribution of the resulting estimators and obtain its estimated variance covariance matrix.


8. Consider a location family $f_\theta(x)=f(x-\theta), \theta \in \Re$, where $f$ is a given density which is symmetric around 0 . Consider estimating $\theta$ by solving an estimating equation of the form

$$
\sum_{i=1}^n \psi\left(X_i-\theta\right)=0,
$$
(a) What should we assume about $\psi$ so that the above estimating equation will lead to an unbiased estimator?

(b) Consider two such estimators $\widehat{\theta}_1$ and $\widehat{\theta}_2$ corresponding to two different $\psi$, say, $\psi_1$ and $\psi_2$. Find them joint asymptotic distribution of $\left(\widehat{\theta}_1, \widehat{\theta}_2\right)$.
[Hint: Work with their influence functions]

9. Prove or disprove (with a counter example): If $\hat{\theta}_1$ and $\hat{\theta}_2$ are asymptotically normal with influence functions $A_1^C\left(X_i\right)$ and $A_2^C\left(X_i\right)$, respectively, then $\widehat{\theta}_1 \widehat{\theta}_2$ is asymptotically normal with IF $\theta_1 A_2^c\left(X_i\right)+\theta_2 A_1^c\left(X_i\right)$

10. Let for $0<p<1, \xi_p=F^{-1}(p)$, be the population $p^{t h}$ quantile, where we assume that $F$ is differentiable at $\xi_p$ with $F^{\prime}\left(\xi_p\right)>0$. In fact assume that $F^{\prime}=f$, where $f$ is the (Lebesgue) density of $F$. As stated in class,
$$
\widehat{\xi}_p-\xi_p=-\frac{F_n\left(\xi_p\right)-p}{f\left(\xi_p\right)}+o_p\left(\frac{1}{\sqrt{n}}\right),
$$
where $\widehat{\xi}_p=\inf \left\{t: F_n(t) \geq p\right\}$ is the sample $\mathrm{p}^{t h}$ quantile.

(a) Show that
$$
\sqrt{n}\left\{\widehat{\sigma}_n-s\right\}=o_p(1),
$$
where
$$
\widehat{\sigma}_n=\sqrt{\frac{1}{n} \sum\left(X_i-\bar{X}\right)^2} \text { and } s=\sqrt{\frac{1}{n-1} \sum\left(X_i-\bar{X}\right)^2}
$$
(b) Work out the joint asymptotic distribution of $\left(s, \widehat{\xi}_{3 / 4}-\widehat{\xi}_{1 / 4}\right)$.

(c) Work on your answer when $F=N\left(\mu, \sigma^2\right)$.

11. Suppose $\left(X_i, Y_i\right)$ are i.i.d. pairs from a bivariate normal distribution such that $\operatorname{corr}\left(X_i, Y_i\right)=\rho$, Find the asymptotic distribution of $\tanh ^{-1}(\widehat{\rho})$, where
$$
\widehat{\rho}=\frac{\sum_{i=1}^n\left(X_i-\bar{X}\right) Y_i}{\sqrt{\left(\sum_{i=1}^n\left(X_i-\bar{X}\right)^2\right)\left(\sum_{i=1}^n\left(Y_i-\bar{Y}\right)^2\right)}} .
$$
How to get $\tanh^{-1} (\hat\rho)$

12. Consider solving the likelihood score equation $S_n(\theta)=0$. Suppose $\widehat{\theta}_c$ is any consistent estimator. The 1-step MLE is defined as
$$
\widehat{\theta}_{1, M L E}=\widehat{\theta}_c-\frac{S_n\left(\widehat{\theta}_c\right)}{S_n^{\prime}\left(\widehat{\theta}_c\right)}
$$
Show that $\widehat{\theta}_{1, M L E} \sim \operatorname{AN}\left(\theta_0, n^{-1} I^{-1}\left(\theta_0\right)\right)$. In other words, the 1 -step MLE has the same asymptotic distribution as the MLE.


13.Consider a simple linear regression model with zero intercept: $Y_i=\beta x_i+\epsilon_i$, where $x_i$ are observed covariates and $\epsilon_i$ are i.i.d. error terms. Let's assume that $\epsilon_i$ have a symmetric distribution.

Consider estimating the regression parameter $\beta$ using an estimating equation of the form
$$
\sum_{i=1}^n \psi\left(Y_i-\beta x_i\right) x_i=0 .
$$
(i) Show that the estimating function used here is unbiased (for zero) under a simple condition on $\psi$ (state the condition).

(ii) Find the asymptotic distribution of $\widehat{\beta}$ including variance estimation.

(iii) Specialize to the case the of the least squares estiumator of $\beta$.

14. Consider a location family $f_\theta(x)=f(x-\theta), \theta \in \Re$, where $f$ is a given density which is symmetric around 0 . Consider estimating $\theta$ by solving an estimating equation of the form
$$
\sum_{i=1}^n \mathrm{w}_i \psi\left(X_i-\theta\right)=0
$$
where $w_i$ are known weights. What should we assume about $\psi$ so that the above estimating equation will lead to an unbiased estimator? Find an expression for an asymptotic confidence interval for $\theta$ based on such an estimator $\widehat{\theta}$.

15. Consider a linear regression model $Y_i=\alpha+\beta X_i+\epsilon_i$, where $X_i$ and $\epsilon_i$ are i.i.d., $\left\{\epsilon_i\right\}$ and $\left\{X_i\right\}$ are independent. In addition assume that $\epsilon_i$ has a finite second moment and a symmetric density. Consider estimating $\theta=(\alpha, \beta)$ using the following estimating equation
$$
\sum_{i=1}^n \psi\left(Y_i-\alpha-\beta X_i\right) \boldsymbol{Z}_i=\mathbf{0},
$$
where $\boldsymbol{Z}_i=\left(1, X_i\right)^T$.

(a) Under what condition on $\psi$, will this estimating equation be unbiased?

(b) Write down the corresponding asymptotic distribution of $\widehat{\theta}$.

(c) Construct an estimator of the asymptotic variance-covariance matrix of $\widehat{\theta}$.

(d) How will you construct a joint confidence set for $(\alpha, \beta)$ using the large sample distribution?

16. Consider data from a simple linear regression without an intercept term, $Y_i=\beta x_i+\epsilon_i$, where $x_i$ are known design constants and $\epsilon_i \sim N(0,1)$ are error terms. Note we are assuming that the error variance is known.

a. Write down the estimating equation for the MLE of the underlying parameter $\beta$.

b. What is the asymptotic distribution of the MLE of $\beta$ ? What condition do you need on the $\left\{x_i\right\}$ for your answer? How will you estimate the asymptotic variance?

c. Find the asymptotic distribution of an $M$-estimator of $\beta$ obtained from the following estimating equation:
$$
\sum_{i=1}^n x_i \psi\left(Y_i-\beta x_i\right)=0
$$
where $\psi$ is an odd bounded function so that $E\left[\psi\left(\epsilon_i\right)\right]=0$.

d. How will your answer change if $\epsilon_i$ are double exponentially distributed?

17. Suppose $X-\mu$ and $-(X-\mu)$ have the same distributions. Obtain the asymptotic distribution of $\widehat{\mu}=\left(\bar{X}+\widehat{\xi}_{0.5}\right) / 2$, where $\widehat{\xi}_{0.5}$ is the sample median.


18. Consider IID data $X_i, 1 \leq i \leq n$, with density
$$
f_\theta(x)=\frac{1}{2 \theta} e^{-|x / \theta|}, x \in \Re, \theta>0 .
$$
Consider estimating $\theta$ by solving an estimating equation
$$
\sum_{i=1}^n\left\{\psi\left(X_i / \theta\right)-1\right\}=0
$$
(i) What condition do you need $\psi$ to satisfy for consistency? 

(ii) Find the asymptotic distribution of the corresponding $\widehat{\theta}$. Simplify your answer as much as possible.

19. Suppose $\widehat{\theta}$ is consistent solution of an estimating equation of the form
$$
\sum_{i=1}^n s\left(X_i, \theta\right)=0 .
$$
(a) What is the influence function of $\widehat{\theta}$ ? (You don't need to give a detailed proof.)

(b) What is the influence function of $(\widehat{\theta})^2$ as an estimator of $\theta^2$ ?

20. Consider data from a regression model $Y_i= \beta x_i+\epsilon_i$, where $x_i$ is non-random and $\epsilon_i \stackrel{\text { i.i.d. }}{\sim} \operatorname{Normal}\left(0, \sigma^2\right)$.

(a) Find the asymptotic distribution of the MLE of $\hat{\beta}$.

(b) Find the asymptotic distribution of a robust estimator of $\beta$ obtained by solving the estimating equation $\sum_{i=1}^n \psi\left(Y_i-\beta x_i\right) x_i=0$, where $\psi$ is a bounded and symmetric function. Construct its variance estimator.

21. Let $X_i, i \geq 1$, be a sequence of i.i.d. random variable and
$$
s^2=\frac{1}{n-1} \sum_{i=1}^n\left(X_i-\bar{X}\right)^2
$$
be the sample variance based on the first $n$ observations.
Find the asymptotic distribution of $s^2$ in two ways: 

(i) using delta method

(ii) using projection arguments for U-statistics. Show sufficient details.

22. (a) Find the influence function of the sample variance $s^2=(n-1)^{-1} \sum_{i=1}^n\left(X_i-\bar{X}\right)^2$ in two ways: 

(i) using the $\Delta$ method, 

(ii) using $U$-statistic techniques.

(b) What is the asymptotic variance of $\sqrt{n} s^2$ ?

23. Consider response $Y_i=\alpha+\beta x_i+\epsilon_i$, where $x_i$ are non-random and $\epsilon_i$ are i.i.d. with $E\left(\epsilon_i\right)=0$.

(a) Write down the least squares estimating equations.

(b) Establish the asymptotic joint normality of your estimators.

(c) Describe how you perform a bootstrap resampling in this situation.

24. (a) Write down the score equations for a logistic regression based on independent samples. 

(b) Obtain the asymptotic distribution of the MLE of the regression parameters and 

(c) an asymptotic $95 \%$ confidence interval for each of the regression parameters.
Calculate (c) for the following artificial data example.

\begin{tabular}{|l|l|l|}
\hline Response & Covariate $Z_1$ & Covariate $Z_2$ \\
\hline 0 & 3 & 1.73 \\
\hline 0 & 3 & 2.54 \\
\hline 1 & 2 & 2.20 \\
\hline 0 & 5 & 1.31 \\
\hline 1 & -1 & 5.40 \\
\hline 0 & 4 & -2.61 \\
\hline 0 & 7 & -1.10 \\
\hline 1 & -5 & 2.12 \\
\hline 1 & 0 & 1.85 \\
\hline
\end{tabular}

25. Let $X_1, \cdots, X_n$ be i.i.d., with density $f_\theta(x)=\theta^{-1} \exp \{-x / \theta\}, x>0$. Under what condition on the function $\gamma$, the following estimating function $\psi(X, \theta)=\gamma(X / \theta)-1$ be unbiased for zero? Give an example of such a $\gamma$. Write down the asymptotic distribution of $\widehat{\theta}$ obtained from the estimating equation
$$
\sum_{i=1}^n\{\gamma(X / \theta)-1\}=0
$$
for a general $\gamma$ satisfying the condition.

26. Consider a simple linear regression model with zero intercept: $Y_i=\beta x_i+\epsilon_i$, where $x_i$ are observed covariates and $\epsilon_i$ are i.i.d. error terms. Let's assume that $\epsilon_i$ have a symmetric distribution.

Consider estimating the regression parameter $\beta$ using an estimating equation of the form
$$
\sum_{i=1}^n \psi\left(Y_i-\beta x_i\right) x_i=0 .
$$
(i) What condition do you impose on $\psi$ so that this equation has a consistent root?

(ii) Show that the least squares estimator is a special case of this.

(iii) Find the asymptotic distribution of the corresponding estimator $\widehat{\beta}$ for a general $\psi$.

(iv) Compute the asymptotic variance (time n ) for three choices of the $\psi$ function and compare.

(v) What happens if the distribution of $\epsilon_i$ is $\mathrm{N}(0,1)$ ? Does it make sense? Also could you have obtained the asymptotic distribution of $\widehat{\beta}$ when $\psi(x)=x$, using any other result you know?

27. Let for $0<p<1, \xi_p=F^{-1}(p)$, be the population $p^{t h}$ quantile, where we assume that $F$ is differentiable at $\xi_p$ with $F^{\prime}\left(\xi_p\right)>0$. In fact assume that $F^{\prime}=f$, where $f$ is the (Lebesgue) density of $F$. As stated in class,
$$
\widehat{\xi}_p-\xi_p=-\frac{F_n\left(\xi_p\right)-p}{f\left(\xi_p\right)}+o_p\left(\frac{1}{\sqrt{n}}\right),
$$
where $\widehat{\xi}_p=\inf \left\{t: F_n(t) \geq p\right\}$ is the sample $\mathrm{p}^{t h}$ quantile.
Work out the asymptotic distribution of $H / I$, where $I:=\widehat{\xi}_{3 / 4}-\widehat{\xi}_{1 / 4}$ and $2 H=\widehat{\xi}_{3 / 4}+\widehat{\xi}_{1 / 4}$ when $F$ is $N(1,1)$. Note that $H / I$ is a robust measure of the coefficient of variation.

28.
$$
\begin{aligned}
&\text { Establish the joint asymptotic normality of }\\
&\left(\frac{\hat{\xi}_{1 / 4}+\hat{\xi}_{3 / 4}}{2}, \quad \frac{\hat{\xi}_{1 / 4}-\hat{\xi}_{3 / 4}}{2}\right), \quad \text { where } \hat{\xi}_p=\inf \left\{x: F_n(x) \geq p\right\}, 0<p<1 .
\end{aligned}
$$

29. Consider $X_1, \ldots, X_n$ i.i.d. from a distribution $F=F_0$ with a location parameter $\theta$.

(a) What does " $F=F_0$ with a location parameter $\theta$ " mean, mathematically?

(b) Suppose $\hat{\theta}_1=\xi_{1 / 2}$ is the sample median and $\hat{\theta}_2$ is another estimator of the location parameter obtained by solving
$$
\sum_{i=1}^n \psi\left(X_i-\hat{\theta}_2\right)=0 .
$$
Find the joint asymptotic distribution of $\left(\hat{\theta}_1, \hat{\theta}_2\right)$ under appropriate regularity conditions (state them as you see fit).

30. (2026 HW3) 

(a) Find the limiting distribution of the sample co-efficient of variation
$$
\widehat{\theta}=\frac{s}{\bar{X}} .
$$

(b) What is its influence function?

31. (2026 HW3) Find the joint asymptotic distribution of the sample mean and the sample median.

32. (2026 HW3) Find the joint asymptotic distribution of the sample first quartile and the sample third quartile.

33. (2026 HW3) Consider i.i.d. sample $X_1, \cdots, X_n$ from $N(\theta, 1)$. Suppose $\widehat{\theta}_1$ is an $M$-estimator obtained as a solution of the estimating equation
$$
\Sigma_{i=1}^n \psi\left(X_i-\theta\right)=0,
$$
with $\psi(u)=\tanh (u)$. Suppose $\widehat{\theta}_2$ is the MLE.
Find out the asymptotic joint distribution of ($\widehat{\theta}_1, \widehat{\theta}_2$) and obtain an estimator of the asymptotic variance-covariance matrix.

34. (2026 Exam 1) Consider binatry data (response) $X_i$ with associated covariate $Z_i, i \leq i \leq n$. We posit a logistic regression model of the form
$$
\operatorname{logit}\left(X_i=1\right)=\beta Z_i / i, 1 \leq i \leq n \text {. }
$$
(a) Write down the likelihood function for the parameter $\beta$.

(b) Write down the likelihood-based score equation.

(c) Obtain the asymptotic distribution of $\widehat{\beta}$, a consistent solution of the score equation (detailed proofs are not needed; however you should identify the quantities in the asymptotic distribution)

(d) Construct a 95\% large sample confidence interval for $\beta$.

35. (2026 Final) Consider a regression model for your response
$$
Y_i=\alpha+\beta x_i+\epsilon_{i},\quad 1 \leq i \leq n,
$$
where $x_i$ are non-random design co-variates. Consider estimating the regression parameters $\theta=(\alpha, \beta)$ using a bivariate estimating equation
where
$$
\begin{aligned}
& S_n(\theta)_{2 \times 1}=0_{2 \times 1} \\
& S_n(\theta):=\sum_{i=1}^n s_i(\theta)_{2 \times 1}=\sum_{i=1}^n\binom{\psi\left(Y_i-\alpha-\beta x_i\right)}{\psi\left(Y_i-\alpha-\beta x_i\right) x_i}_{2 \times 1}
\end{aligned}
$$
Suppose $\psi$ is an odd function and $\epsilon_i$ has a symmetric (around 0 ) distribution.

(a) Show that $s_i(\theta)$ has mean 0 (vector), under the true $\theta$.

(b) Using the arguments/results provided in class obtain the asymptotic distribution of $\widehat{\theta}$.

Simplify your answer as much as possible. Write down any condition you need on the design values $x_i$.

(c) Estimate the asymptotic variance-covariance matrix of $\hat{\theta}$.

\newpage

\section{U-statistics}
1. Consider i.i.d. $X_1, \cdots, X_n$ from a common distribution $F$, with $\int x^2 d F<\infty$.
(a) Show that
$$
\widehat{\sigma}_n^2:=\frac{1}{n} \sum_{i=1}^n\left(X_i-\bar{X}\right)^2
$$
is a $V$-statistics

(b) Using the projection result for $U / V$ statistics, obtain the asymptotic distribution of $\widehat{\sigma}_n^2$

2. Suppose we have binary response $Y_i$ (0 or 1) and a single cavariate $Z_i, 1 \leq i \leq n$. Suppose, we fit a logistic regression model of the form
$$
\operatorname{logit}\left(\operatorname{Pr}\left\{Y_i=1 \mid Z_i\right\}\right)=\beta \mathrm{Z}_i, 1 \leq i \leq n .
$$
Write down the score test for $H_0: \beta=0$ against $H_1: \beta \neq 0$.

3. Find a kernel $h$ that $E h=\mu_4$, the forth central moment. Find the asymptotic distribution of the corresponding U-statistic and write down the formula for an approximate (large sample) confidence interval for $\mu_4$.

4. Suppose $X_i, 1 \leq i \leq n_1$, and $Y_j, 1 \leq j \leq n_2$ sample from two independent populations with CDFs, $F$ and $G$, respectively. Compute the asymptotic distribution of the test statistic based on a two-sample U-statistic with kernel
$$
h\left(X_i, Y_j\right)=I\left(Y_j>X_i\right),
$$
under the null hypothesis $H_0: F=G$.
What is the name of this test?

5. Consider samples $X_1, \cdots, X_{n_1}$ and $Y_1, \cdots, Y_{n_2}$ from two independent populations. Consider the two-sample U-statistic corresponding to the kernel
$$
h\left(X_1, \cdots, X_{m_1} ; Y_1, \cdots, Y_{m_2}\right)=h^1\left(X_1, \cdots, X_{m_1}\right) h^2\left(Y_1, \cdots, Y_{m_2}\right) .
$$
where $h^1$ and $h^2$ are kernels of orders $m_1$ and $m_2$, respectively. Find its asymptotic distribution.

6. Consider i.i.d. $X_1, \cdots, X_n$ from a common distribution with mean $\mu$ and variance $\sigma^2$. Let $V$ be a V-statistic for estimating $\mu^2$ and $U$ be a U-statistic for estimating $\sigma^2$.

(i) Write down $V$ and $U$ in terms of $X_1, \cdots, X_n$. Are they unbiased for the respective parameters?

(ii) Obtain the joint asymptotic distribution of $(V, U)$.

(iii) Discuss how you will obtain an approximate joint confidence set for $\left(\mu^2, \sigma^2\right)$.\


7. Find the U-statistics $U_1$ and $U_2$ that are unbiased for $\sigma^2$ and $\mu_3$, respectively. Obtain the joint asymptotic distribution of $\left(U_1, U_2\right)$.

8. Consider two independent samples $\left\{X_1, \cdots, X_n\right\}$ and $\left\{Y_1, \cdots, Y_m\right\}$ distributed $F$ and $G$, respectively (assume both $F$ and $G$ are continuous). Work out the asymptotic null distribution of the test statistic
$$
\frac{1}{n m} \sum_{i=1}^n \sum_{j=1}^m I\left(X_i<Y_j\right)
$$
under $H_0: F=G$.

9. Suppose $h^{1}$ and $h^{2}$ are two kernels of orders $p$ and $q$, respectively. Consider two independent samples $\left\{X_{1}, \cdots, X_{n}\right\}$ and $\left\{Y_{1}, \cdots, Y_{m}\right\}$ distributed $F$ and $G$, respectively.
$$
\text { Let } E_{F} h^{1}\left(X_{1}, \cdots, X_{p}\right)=\eta \text { and } E_{\mathrm{G}} h^{2}\left(Y_{1}, \cdots, Y_{q}\right)=\kappa \text {. }
$$
(i) Obtain the asymptotic distribution of a two-sample $U$ statistic with kernel
$$
\mathcal{H}\left(X_{1}, \cdots, X_{p} ; Y_{1}, \cdots, Y_{q}\right)=\sqrt{h^{1}\left(X_{1}, \cdots, X_{p}\right) \cdot h^{2}\left(Y_{1}, \cdots, Y_{q}\right)}
$$
(ii) Suppose $\eta=\sigma_{F}^{2}$ and $\kappa=\sigma_{G}^{2}$. Find a consistent estimator of $\sigma_{F} \sigma_{G}$ and obtain its asymptotic distribution.

10. Consider two independent, i.i.d. samples $\left\{X_{11}, \cdots, X_{1 n_1}\right\}$ and $\left\{X_{21}, \cdots\right.$, $\left.X_{2 n_2}\right\}$ from two populations. Let $U_1$ and $U_2$ be two U-statistics based on each of these samples with kernels $h^{(1)}$ and $h^{(2)}$, respectively. Find the large sample distribution of the statistic $T=g\left(U_1, U_2\right)$, where $g$ is a smooth function. Show the detailed steps and write down the condition on $n_1$ and $n_2$ for the result.

11. Suppose $T_1$ is a U-statistic based on a kernel $h_1$ and $T_2$ is a V-statistic with kernel $h_2$.

(a) What is the influence function of the ratio $T_2 / T_1$?

(b) Using (2a), obtain the asymptotic normality of
$$
\frac{\bar{X}^2}{s^2}, \quad \text { where } s^2=(n-1)^{-1} \sum_{i=1}^n\left(X_i-\bar{X}\right)^2 .
$$

12. Find the influence function of the sample variance
$$
s^2=(n-1)^{-1} \sum_{i=1}^n\left(X_i-\bar{X}\right)^2
$$
in two ways:
(i) using the $\Delta$-method,
(ii) using U-statistic techniques.
(b) What is the asymptotic variance of $\sqrt{n} s^2$?

13. Find two kernels $h_1$ and $h_2$ which are unbiased for the population variance $\sigma^2$ and the third central moment $\mu_3$, respectively. Let $U_1$ and $U_2$ denote the corresponding U-statistics.

(a) Obtain the asymptotic distribution of
$$
\frac{U_2}{\sqrt{U_1}} .
$$
(b) Construct a large-sample confidence interval for $\mu_3 / \sigma$.

14. Suppose
$$
U_1=U_1\left(X_{11}, \ldots, X_{1 n_1}\right) \quad \text { and } \quad U_2=U_2\left(X_{21}, \ldots, X_{2 n_2}\right)
$$
are U-statistics of orders $m_1$ and $m_2$, respectively, based on independent samples $\left\{X_{11}, \ldots, X_{1 n_1}\right\}$ and $\left\{X_{21}, \ldots, X_{2 n_2}\right\}$. Suppose $\theta_2=\mathbb{E}\left[U_2\left(X_{21}, \ldots, X_{2 n_2}\right)\right] \neq 0$. Consider the two-sample statistic
$$
R=\frac{U_1}{U_2} .
$$
(a) Find the asymptotic distribution of $R$ and indicate how you will estimate its asymptotic variance (no jackknife or bootstrap please)

(b) Apply the general result in (a) to obtain the asymptotic distribution of
$$
\frac{s_1^2}{s_2^2}, \quad \text { where } s_k^2 \text { is the sample variance of the } k \text { th sample, } k=1,2 \text {. }
$$
(c) Write down an asymptotic linear representation for $R$.

15. Consider a one-sample U-statistic
$$
U=\frac{2}{n(n-1)} \sum_{1 \leq i<j \leq n} I\left(X_i+X_j>0\right) .
$$
Find the asymptotic distribution of $U$ under the assumption that the $X_i$ have a symmetric distribution (i.e. $X_i \stackrel{d}{=}-X_i$ ).

16. Consider i.i.d. bivariate pairs $Z_i=\left(X_i, Y_i\right), 1 \leq i \leq n$. Consider the $U$-statistic
$$
U=\frac{1}{\binom{n}{2}} \sum_{i<j} \operatorname{sgn}\left(X_i-X_j\right) \operatorname{sgn}\left(Y_i-Y_j\right),
$$
where $\operatorname{sgn}(x)=-1,0,1$, for $x<0,=0,>0$.

(i) What is the name of this U-statistic?

(ii) Find the asymptotic distribution of $U$, under the null hypothesis $H_0: X_i$ and $Y_i$ are independent.

17. Suppose $X_i$ 's are i.i.d, $1 \leq i \leq n$,
$$
V_n=\frac{1}{n^m} \sum_{i_1=1}^n \sum_{i_2=1}^n \cdots \sum_{i_m=1}^n h\left(X_{i_1}, \cdots, X_{i_m}\right)
$$
be a $V$-statistic based on a symmetric and positive kernel $h$, and
$$
\theta=E h\left(X_1, \cdots, X_m\right), 0<E h^2\left(X_1, \cdots, X_m\right)<\infty .
$$
Let $g:(0, \infty) \rightarrow(0, \infty)$ be a smooth function.

(a) What will be the influence function of $g\left(V_n\right)$ as an estimator of $g(\theta)$ ? Give arguments for your answer.

(b) Apply your answer in part (a) to find the influence function $I F\left(X_i\right)$ of
$$
\widehat{\sigma}_n=\sqrt{\frac{1}{n} \sum_{i=1}^n\left(X_i-\bar{X}\right)^2}
$$
(c) How will you estimate the population quantities appearing in $I F\left(X_i\right)$ to obtain an estimated influence function, say, $\widehat{F}\left(X_i\right)$?

(d) Using $\widehat{F}\left(X_i\right)$ 's, construct a large sample $95 \%$ confidence interval for $\sigma$.

18. (2026 Exam 2) Consider IID binary data $X_1, \cdots, X_n$ with $p=\operatorname{Pr}\left\{X_1=0\right\}$.
A natural estimator of $p$ is $\widehat{p}=n^{-1} \sum_{i=1}^n I\left(X_i=0\right)$ and its asymptotic variance is $p(1-p) / n$. 

Express $\widehat{p}(1-\widehat{p})$ as a $U$ or $V$ statistic (show details) and obtain its asymptotic distribution.

19. (2026 Exam 2) Suppose we have independent samples 
$$
\left\{X_{11}, \cdots, X_{1 n_1}\right\},\left\{X_{21}, \cdots, X_{2 n_2}\right\},\left\{X_{31}, \cdots, X_{3 n_3}\right\}
$$ 
from 3 populations with means $\mu_1, \mu_2, \mu_3$, respectively.

A natural estimator of $\theta=\mu_1 \times \mu_2 \times \mu_3$ is $\widehat{\theta}:=\bar{X}_1 \times \bar{X}_2 \times \bar{X}_3$,, the product of the 3 sample means.

(a) $\widehat{\theta}$ is both a multi-sample $U$ and a multi-sample $V$ statistic. True or False?

(b) Obtain the asymptotic distribution of $\widehat{\theta}$. Try to use the notations used in class and show the necessary steps/calculations. Also, provide the conditions on the sample sizes for the normal limit to hold.

20. (2026 HW5) Find U-statistics $U_1$ and $U_2$ that are unbiased for $\mu$ and $\sigma^2$, respectively. Obtain the asymptotic distribution of $\sqrt{U_2} / U_1$. What is this quantity called in statistics?

21. (2026 HW5) Suppose $h^1$ and $h^2$ are two kernels of orders $p$ and $q$, respectively. Consider two independent samples $\left\{X_1, \cdots, X_n\right\}$ and $\left\{Y_1, \cdots, Y_m\right\}$ distributed $F$ and $G$, respectively.
$$
\text { Let } E_F h^1\left(X_1, \cdots, X_p\right)=\eta \text { and } E_{\mathrm{G}} h^2\left(Y_1, \cdots, Y_q\right)=\kappa \text {. }
$$
(i) Obtain the asymptotic distribution of a two-sample $U$ statistic with kernel
$$
\mathcal{H}\left(X_1, \cdots, X_p ; Y_1, \cdots, Y_q\right)=\sqrt{h^1\left(X_1, \cdots, X_p\right) \cdot h^2\left(Y_1, \cdots, Y_q\right)} .
$$
(ii) Suppose $\eta=\sigma_F^2$ and $\kappa=\sigma_G^2$. Find a consistent estimator of $\sigma_F \sigma_G$ and obtain its asymptotic distribution.

22. (2026 HW5). Consider two independent samples $\left\{X_1, \cdots, X_n\right\}$ and $\left\{Y_1, \cdots, Y_m\right\}$ distributed $F$ and $G$, respectively (assume both $F$ and $G$ are continuous). Work out the asymptotic null distribution of the test statistic

$$
\frac{1}{n m} \sum_{i=1}^n \sum_{j=1}^m I\left(X_i<Y_j\right)
$$

under $H_0: F=G$. What is the name of this test?

23. (2026 Final) (a) Consider a 2-sample U-statistic with a general kernel $h$ satistying the usual symmetry and finite second moment assumptions
$$
U=\frac{1}{\binom{n_1}{m_1}\binom{n_2}{m_2}} \sum_{1 \leq i_1 \cdots<i_{m_1} \leq n_1} \sum_{1 \leq j_1 \cdots<j_{m_2} \leq n_2} h\left(X_{1 i_1}, \cdots, X_{1 i_{m_1}} ; X_{2 j_1}, \cdots, X_{2 j_{m_2}}\right) .
$$
Suppose $\min \left(n_1, n_2\right) \rightarrow \infty$, such that $n_1 / N \rightarrow \alpha, N=n_1+n_2$.
Assuming the following linear approximation
$$
U^c=\frac{m_1}{n_1} \sum_{i=1}^{n_1} h_{11}^c\left(X_{1 i}\right)+\frac{m_2}{n_2} \sum_{j=1}^{n_2} h_{21}^c\left(X_{2 j}\right)+o_p\left(N^{-1 / 2}\right),
$$
find the asymptotic distribution of $U$.

(b) Work out your answer in (a) for the statistic
$$
\frac{1}{n_1 n_2} \sum_{i=1}^{n_1} \sum_{j=1}^{n_2} I\left(X_{1 i}+X_{2 j}>0\right),
$$
expressing your calculations in terms of $F$ and $G$, the CDFs of $X_{1 i}$ and $X_{2 j}$, respectively.

24. (2025 Final) Suppose $X_i$ 's are i.i.d, $1 \leq i \leq n$,
$$
V_n=\frac{1}{n^m} \sum_{i_1=1}^n \sum_{i_2=1}^n \cdots \sum_{i_m=1}^n h\left(X_{i_1}, \cdots, X_{i_m}\right)
$$
be a $V$-statistic based on a symmetric and positive kernel $h$, and
$$
\theta=E h\left(X_1, \cdots, X_m\right), 0<E h^2\left(X_1, \cdots, X_m\right)<\infty .
$$
Let $g:(0, \infty) \rightarrow(0, \infty)$ be a smooth function.

(a) What will be the influence function of $g\left(V_n\right)$ as an estimator of $g(\theta)$ ? Give arguments for your answer.

(b) Apply your answer in part (a) to find the influence function IF $\left(X_i\right)$ of
$$
\widehat{\sigma}_n=\sqrt{\frac{1}{n} \sum_{i=1}^n\left(X_i-\bar{X}\right)^2}
$$

(c) How will you estimate the population quantities appearing in $I F\left(X_i\right)$ to obtain an estimated influence function, say, $\widehat{F}\left(X_i\right)$?

(d) Using $\widehat{F}\left(X_i\right)$ 's, construct a large sample $95 \%$ confidence interval for $\sigma$.

\newpage
\section{Jackknife and Bootstrap}
1. Consider a random sample $\{1,3,4\}$ from a population. Suppose we want to estimate the population median by the sample median (use any acceptable definition; just write down which version you are using).

(a) Compute the jackknife pseudo values.

(b) Compute a variance estimate of the sample median using the jackknife.

2. Generate data with various sample sizes $(n=30,200,500)$ from two different distributions. In each setting, simulate 1000 data sets. For each data set, obtain the following

(ii) Estimated variance of $\widehat{\xi}_{1 / 2}$ using the jackknife.

3. Consider i.i.d. data $X_1, \cdots, X_n$ and the following robust measure of the coefficient of variation
$$
\widehat{\theta}=\frac{\widehat{\tau}}{M},
$$
where
$$
\widehat{\tau}=\frac{1}{n^2} \sum_{i=1}^n \sum_{j=1}^n\left|X_i-X_j\right|
$$
and $M$ is the sample median.
Construct a variance estimator for $\widehat{\theta}$ using each of the following methods

(a) the large sample theory techniques (influence function delta method etc)

(b) the jackknife

(c) the bootstrap

Which one is not likely to work?

4. Consider independent samples $X_{11}, \cdots, X_{1 n_1}$ and $X_{21}, \cdots, X_{2 n_2}$ from two populations $F$ and $G$. Let $\widehat{\theta}_{n_1, n_2}=\widehat{\theta}\left(X_{11}, \cdots, X_{1 n_1} ; X_{21}, \cdots, X_{2 n_2}\right)$ be an estimator of some population parameter $\theta=\theta(F, G)$. Propose a Jackknife bias corrected version of $\widehat{\theta}$.


5. Suppose $n=3$ and $\widehat{\theta}_{-1}=-1, \widehat{\theta}_{-2}=0, \widehat{\theta}_{-3}=1$. Computh the jackknife estimate of the variance of $\widehat{\theta}$ using Efron's formula.


6. How should we use jackknife in a two sample problem, such as with a two sample U-statistic? Write down your thoughts

7. Write down two formulas for the jackknife estimator of variance of a statistic and prove their equality.

8. Suppose $\theta$ is an estimator based on i.i.d. data $X_1, \cdots, X_n$. Describe how vou will construct a bootstrap estimator of the density of $\widehat{\theta}$.

9. Consider a coin toss experiment leading to $X=0$ or 1 , according to whether the coin shows a tail or a head. Suppose $n=3, X_1=1, X_2=0$, and $X_3=0$. Consider the standard bootstrap resampling procedure. Obtain the following:

(i) $P^*\left\{X_2^*=1\right\}$, and

(ii) $E^*\left(\overline{X^*}\right)$.

Note: These will be exact answers. Do not give a Monte Carlo algorithm.

10. Consider a $V$-statistic $T$ based on a kernel $h$ of degree $m$.

(a) Write down the asymptotic linear representation for $T$.

(b) Suppose the remainder term in this approximation can be ignored. Obtain an estimate for the variance of $T$ using the jackknife principle of the main term.

11. (2026 Final) Consider a sample of size 5 . Suppose the following values were obtained:
$$
\widehat{\theta}_{-1}=3, \widehat{\theta}_{-2}=3, \widehat{\theta}_{-3}=0, \widehat{\theta}_{-4}=1, \widehat{\theta}_{-5}=1 .
$$
Compute the jacknife estimate of the variance of $\widehat{\theta}$ for the same sample.
\end{document}
